\section{Reed--Solomon Projection and Proximity}\label{sec:rs-proximity-projection}

This section develops the HCP-specific application bridge. The companion
coding paper \cite{vanDerVlugt2026proximity} treats universal list decoding
and pair MCA over the full ambient word space. We retain the coding
definitions needed to state the bridge and its cited endpoint bound here;
the companion's capacity-grid and large-field results are separate from this
section's fixed-field instantiation.

The staged projection maps are
\[
  \operatorname{Pl}(\mathcal F)
  \xrightarrow{\ E\ }
  c
  \xrightarrow{\ \operatorname{proj}_{\mathcal E}\ }
  p_c
  \xrightarrow{\ \Phi_\ast\ }
  y.
\]
Here \(c\) is the chain's cochain payload, \(p_c\) is its plane shadow, and
\(y\) is the finite-field received word. The public-binding relation introduced
below identifies \(c\) with an accepted cochain commitment.
For a valid frame, the corresponding stage predicates are
\[
  \sigma(c)\in Z,
  \qquad
  p_c\in\mathcal C_{\mathcal P},
  \qquad
  y\in T.
\]
Thus the local syndrome and canonical Reed--Solomon target are acceptance
conditions on projected objects, not additional projection maps. The chain has
two complementary proof directions. The stage maps give codeword preservation
and the near-side metric handoff
\[
  D(y)\le d_L(c)\le \rho(\mathcal F),
\]
while a supplied final far-region theorem gives the calibrated non-collapse
implication needed for a gap reduction. Once both directions have been proved,
they yield the end-to-end proximity statement. The projective Pl\"ucker
coordinate is the geometric handoff: a private frame determines
\(\operatorname{Pl}(\mathcal F)\), and the encoder \(E\) produces the chain
cochain \(c\). At the public boundary, the relational form \(\mathcal R\)
identifies that payload with an accepted cochain commitment.

At the projection-chain level, the private source datum is the frame witness
\(\mathcal F\). When this chain is integrated with the semantic and commitment
layers of \cref{sec:mathematical-framework,sec:public-artifacts-commitments}, that
frame may be accompanied by fibration data, paths, and commitment openings.
The verifier-facing mathematical data then consist of the accepted cochain
commitment, the Pl\"ucker coordinate bound to it, the plane projection, and the
finite received word. Hashes and roots are compressed handles for canonical
bytes of those data, not replacements for the mathematical objects.

\subsection{The End-to-End Projection Contract}

The displayed chain must do more than produce a finite word.  On its near
side, it must carry every valid frame to a Reed--Solomon codeword and control
the target distance by the frame disagreement:
\[
  \mathcal F\in\mathcal V_{\mathrm{fr}}
  \Longrightarrow y\in T,
  \qquad
  D(y)\le \rho(\mathcal F).
\]
On its far side, it must not collapse a source instance beyond the calibrated
target radius.  The stage maps prove the first two statements; a separate
far-region witness must prove the non-collapse statement.  These are precisely
the two directions required to pull a target gap solver back through the
composed map.

We now isolate that contract in notation that will be instantiated by the
concrete stages below.  A source proximity datum is the three-field record
\[
  \mathcal X=(X,G_X,\rho_X),
\]
where \(X\) is the source carrier, \(G_X\subseteq X\) is its exact or valid
locus, and \(\rho_X:X\to\mathbb R\) is its disagreement score.  End to end,
these are respectively the frame-witness carrier, the valid-frame locus, and
the Pl\"ucker disagreement.  At an intermediate stage the same interface may be
instantiated by committed cochains and their local syndrome distance.  A
code-side proximity datum is the three-field record
\[
  \mathcal Y=(Y,\mathcal C_Y,d_Y),
\]
where \(Y\) is the target word carrier, \(\mathcal C_Y\subseteq Y\) is
the code, and \(d_Y:Y\to\mathbb R\) is the distance-to-code score.

\begin{definition}[Evaluation projection]\label{def:evaluation-projection}
An evaluation projection \(\eta:\mathcal X\to\mathcal Y\) is a map
\[
  \eta:X\longrightarrow Y
\]
such that
\[
  x\in G_X\Longrightarrow \eta(x)\in\mathcal C_Y,
  \qquad
  d_Y(\eta(x))\le \rho_X(x)
  \quad(x\in X).
\]
Consequently, if \(\rho_X(x)\le \varepsilon\), then
\[
  d_Y(\eta(x))\le \varepsilon.
\]
\end{definition}

A promised gap-proximity problem is a tuple
\[
  P=(X,s,r_{\mathrm{near}},r_{\mathrm{far}})
  \qquad
  \text{with}\qquad r_{\mathrm{near}}\le r_{\mathrm{far}}.
\]
A promised instance has near and far loci
\[
  P_{\mathrm{near}}=\{x\in X:s(x)\le r_{\mathrm{near}}\},
  \qquad
  P_{\mathrm{far}}=\{x\in X:r_{\mathrm{far}}\le s(x)\}.
\]
A solver \(A:X\to\{\operatorname{accept},\operatorname{reject}\}\) is correct
on the promise when
\[
  x\in P_{\mathrm{near}}\Rightarrow A(x)=\operatorname{accept},
  \qquad
  x\in P_{\mathrm{far}}\Rightarrow A(x)=\operatorname{reject}.
\]
These closed promises are disjoint when
\(r_{\mathrm{near}}<r_{\mathrm{far}}\). If the radii coincide at \(r\), a
correct accept/reject solver can exist only when the boundary shell
\(\{x:s(x)=r\}\) is empty: an attained boundary would require both decisions.
The solver pullbacks below are conditional on a correct target solver; they
do not assert the existence of an efficient solver.
The canonical decision consequence in \cref{prop:canonical-bridge} uses a
strictly smaller near radius than the Johnson far radius, while leaving the
list-decoding and MCA radii unchanged.

A gap reduction \(f:P\to Q\) preserves near instances and far instances. Hence
any correct solver for \(Q\) pulls back to a correct solver for \(P\).

\begin{proposition}[Projection and solver pullback]\label{prop:projection-pullback}
Let \(\eta:\mathcal X\to\mathcal Y\) be an evaluation projection.
Then:
\[
  \rho_X(x)\le\varepsilon\Longrightarrow d_Y(\eta(x))\le\varepsilon.
\]
If, in addition, the projection is calibrated by radii
\[
  r^X_{\mathrm{near}}\le r^X_{\mathrm{far}},
  \qquad
  r^Y_{\mathrm{near}}\le r^Y_{\mathrm{far}},
  \qquad
  r^X_{\mathrm{near}}\le r^Y_{\mathrm{near}},
\]
and satisfies
\[
  r^X_{\mathrm{far}}\le\rho_X(x)
  \Longrightarrow
  r^Y_{\mathrm{far}}\le d_Y(\eta(x)),
\]
then \(\eta\) is a gap reduction from the source gap problem to the target gap
problem. Consequently, every correct target solver \(A\) pulls back to the
correct source solver \(A\circ\eta\).
\end{proposition}

\begin{proof}
The first assertion is transitivity:
\[
  d_Y(\eta(x))\le\rho_X(x)\le\varepsilon.
\]
For the calibrated near side, if \(x\) is source-near, then
\[
  d_Y(\eta(x))
  \le \rho_X(x)
  \le r^X_{\mathrm{near}}
  \le r^Y_{\mathrm{near}},
\]
so \(\eta(x)\) is target-near. The far side is exactly the displayed
non-collapse implication. If \(A\) is correct on the target promise, then for
a source-near \(x\), the just-proved near preservation gives
\[
  A(\eta(x))=\operatorname{accept},
\]
and for a source-far \(x\), the far preservation gives
\[
  A(\eta(x))=\operatorname{reject}.
\]
This is precisely correctness of the pulled-back solver.
\end{proof}

For a lossy projection chain, the final interface is calibrated. A calibrated
end-to-end gap projection consists of an evaluation projection
\(\eta:\mathcal X\to\mathcal Y\), radii
\[
  r^X_{\mathrm{near}}\le r^X_{\mathrm{far}},
  \qquad
  r^Y_{\mathrm{near}}\le r^Y_{\mathrm{far}},
\]
the near compatibility
\[
  r^X_{\mathrm{near}}\le r^Y_{\mathrm{near}},
\]
and the far non-collapse implication
\[
  r^X_{\mathrm{far}}\le \rho_X(x)
  \Longrightarrow
  r^Y_{\mathrm{far}}\le d_Y(\eta(x)).
\]
The ordinary same-radius projection is the special case
\(r^X_{\mathrm{near}}=r^Y_{\mathrm{near}}\) and
\(r^X_{\mathrm{far}}=r^Y_{\mathrm{far}}\).

\subsection{Grassmannian Source and Pl\"ucker Quotient}

The first Grassmannian context scores an ambient observation against one fixed
harmonic subspace. Let \(W\le \mathbb R^n\) be a Grassmannian point with
orthogonal projection \(P_W\), and put
\[
  \Delta_W=I-P_W.
\]
Its carrier is \(\mathbb R^n\), the exact locus is
\[
  G_W=W,
\]
and the disagreement score is
\[
  \rho_W(b)=\|P_{W^\perp}b\|.
\]
The residual identity is
\[
  \rho_W(b)=\operatorname{res}_W(b),
\]
and
\[
  \rho_W(b)=0
  \Longleftrightarrow
  b\in W
  \Longleftrightarrow
  b\in\ker(\Delta_W).
\]
The identity map is therefore an evaluation projection from this vector-level
Grassmannian proximity datum to its harmonic code shadow.

\begin{proposition}[Grassmannian residual]\label{prop:grassmannian-residual}
For \(W\le\mathbb R^n\) with orthogonal projection \(P_W\),
\[
  \rho_W(b)=\|P_{W^\perp}b\|
  =
  \operatorname{res}_W(b),
\]
and
\[
  \rho_W(b)=0
  \Longleftrightarrow b\in W
  \Longleftrightarrow b\in\ker(\Delta_W).
\]
\end{proposition}

\begin{proof}
The projection model has
\[
  b=P_Wb+P_{W^\perp}b,\qquad
  \Delta_W b=(I-P_W)b=P_{W^\perp}b.
\]
Thus the Laplacian residual is
\[
  \operatorname{res}_W(b)=\|\Delta_W b\|=\|P_{W^\perp}b\|=\rho_W(b).
\]
Since a norm vanishes exactly on the zero vector,
\[
  \rho_W(b)=0
  \Longleftrightarrow P_{W^\perp}b=0.
\]
The latter condition is equivalent to \(b=P_Wb\), hence to \(b\in W\). In the
projection Hodge model, harmonicity is membership in \(\ker(\Delta_W)=W\).
\end{proof}

The composed chain below uses a second, downstream Grassmannian context: its
source is a frame witness and its disagreement \(\rho(\mathcal F)\) is supplied
with that frame-level proximity instance. The vector residual \(\rho_W(b)\) and
the frame score \(\rho(\mathcal F)\) are therefore not identified merely by
sharing the same underlying harmonic subspace; an application relating them
must state the comparison explicitly.

Now fix a rank \(r\) Grassmannian point. A frame witness is
\[
  \mathcal F=(W;w_1,\ldots,w_r)
\]
with \(w_i\in W\), linearly independent, and
\[
  \operatorname{span}_{\mathbb R}\{w_1,\ldots,w_r\}=W.
\]
Its exterior coordinate is
\[
  x_{\mathcal F}=w_1\wedge\cdots\wedge w_r
  \in \bigwedge^r\mathbb R^n.
\]
The raw Pl\"ucker validity conditions are
\[
  x_{\mathcal F}\ne 0,\qquad
  x_{\mathcal F}\text{ is decomposable},\qquad
  \mathsf{PluckerRelations}(x_{\mathcal F}).
\]
Here \(\mathsf{PluckerRelations}\) is the relation selected by the framed
system; the present unrestricted projective instantiation takes it to be
decomposability. Public-shadow validity additionally requires normalization and
support admissibility. Let \(\mathcal V_{\mathrm{fr}}\) denote the locus of
frame witnesses satisfying both the raw-coordinate and public-shadow validity
conditions.
In the concrete projective instantiation,
\[
  x\sim y
  \Longleftrightarrow
  \exists u\in\mathbb R^\times,\ y=u x,
\]
and the Pl\"ucker coordinate is the projective class
\[
  \operatorname{Pl}(\mathcal F)=[x_{\mathcal F}]
  \in \mathbb P\!\left(\bigwedge^r\mathbb R^n\right).
\]
The soundness conditions are:
\[
\begin{gathered}
  x_{\mathcal F}\ne 0,\qquad
  \mathsf{PluckerRelations}(x_{\mathcal F}),\\
  \operatorname{Pl}(\mathcal F)\text{ is normalized/support-admissible},\\
  \mathcal F.\mathsf{point}=\mathcal G.\mathsf{point}
    \Longrightarrow x_{\mathcal F}\sim x_{\mathcal G},\qquad
  x\sim y\Longrightarrow [x]=[y].
\end{gathered}
\]
Thus changing the private frame for the same Grassmannian point does not change
the projective Pl\"ucker coordinate.

\begin{proposition}[Projective Pl\"ucker invariance]\label{prop:plucker-invariance}
Assume the displayed Pl\"ucker soundness conditions. If two frame witnesses
\(\mathcal F=(W;w_1,\ldots,w_r)\) and
\(\mathcal G=(W;v_1,\ldots,v_r)\) present the same Grassmannian point, then
\[
  [w_1\wedge\cdots\wedge w_r]
  =
  [v_1\wedge\cdots\wedge v_r]
  \in \mathbb P\!\left(\bigwedge^r\mathbb R^n\right).
\]
\end{proposition}

\begin{proof}
Let \(x_{\mathcal F}=w_1\wedge\cdots\wedge w_r\) and
\(x_{\mathcal G}=v_1\wedge\cdots\wedge v_r\). The soundness condition for
frame change gives \(x_{\mathcal F}\sim x_{\mathcal G}\). By definition of the
projective quotient, this means that for some \(u\in\mathbb R^\times\),
\[
  x_{\mathcal G}=u x_{\mathcal F}.
\]
The quotient map to projective space identifies nonzero scalar multiples; hence
\([x_{\mathcal F}]=[x_{\mathcal G}]\).
\end{proof}

For the concrete exterior-power quotient these algebraic soundness inputs are
now proved, rather than left as additional hypotheses. Independent frames have
nonzero exterior product, and two bases of the same rank-\(r\) subspace give
nonzero elements of its one-dimensional top exterior power, hence differ by
a nonzero scalar.

\begin{proposition}[Recovery from the full projective point]
\label{prop:plucker-recovery}
Let \(w_1,\ldots,w_r\) be an independent frame for \(W\subseteq\mathbb R^n\),
and let \(c=w_1\wedge\cdots\wedge w_r\). Then
\[
 W=\{x\in\mathbb R^n:x\wedge c=0\}.
\]
This reconstruction depends only on \([c]\). Consequently, two genuine
rank-\(r\) frame witnesses have the same full projective point if and only
if they present the same subspace. At the fixed ambient inner product, the
full point determines \(P_W\) and \(I-P_W\).
\end{proposition}
\begin{proof}
If \(x\in W\), the frame with \(x\) adjoined is dependent, so its exterior
product vanishes. If \(x\notin W\), the adjoined frame is independent and its
exterior product is nonzero. Multiplication of \(c\) by a nonzero scalar
does not change its annihilator. The projector is then the unique orthogonal
projection onto the reconstructed subspace.
\end{proof}

The Lean construction \lean{Concrete.recoverSubspace} descends this
annihilator through the existing projective quotient;
\lean{recoverHarmonicProjector} reconstructs its orthogonal projector.
The raw quotient also contains the zero orbit, whose annihilator is the whole
ambient space; the recovery proposition is indexed by genuine frame witnesses.
If \(W=\ker\Delta_x\), full projective data together with \(b\) determine
the operative target \(P_Wb\). This is mathematical recovery on exact real
coordinates, with no claim of an efficient finite-encoding algorithm. A
digest-only public value is a different observation and requires its own
binding and visibility analysis. The full point does not determine the
positive spectrum of the original \(\Delta_x\).

The source metric gate states
\[
  0\le \rho(\mathcal F),\qquad
  \operatorname{Pl}(\mathcal F)=\operatorname{Pl}(\mathcal G)
  \Longrightarrow
  \rho(\mathcal F)=\rho(\mathcal G).
\]
The Pl\"ucker metric gate fixes the admissible projective chart and a projective
distance \(d_{\mathrm{Pl}}\) satisfying
\[
  0\le d_{\mathrm{Pl}}(\mathcal F,\mathcal G),
  \qquad
  \operatorname{Pl}(\mathcal F)=\operatorname{Pl}(\mathcal F')
  \Longrightarrow
  d_{\mathrm{Pl}}(\mathcal F,\mathcal G)
  =d_{\mathrm{Pl}}(\mathcal F',\mathcal G).
\]

\subsection{Committed Cochains and Local Syndrome}

The committed-cochain carrier in the proximity bridge is exactly the cochain
commitment payload of \cref{sec:public-artifacts-commitments}. Write
\[
  \mathcal C_{\mathrm{com}}
\]
for that carrier. An element is the payload
\[
  c=(q,n,\mathbf c,\tau)
\]
with supported modulus, declared edge count, edge values, and topology hash.
Let
\[
  \mathcal C_{\mathrm{com}}^{\mathrm{wf}}\subseteq\mathcal C_{\mathrm{com}},
  \qquad
  \mathcal S_{\mathrm{coh}}
\]
denote respectively the well-formed commitment payloads and the coherent sheaf
data. For \(S\in\mathcal S_{\mathrm{coh}}\), define
\[
  G_S=\{c\in\mathcal C_{\mathrm{com}}:
  c\in\mathcal C_{\mathrm{com}}^{\mathrm{wf}},\quad
  c\mathrel{\mathsf{binds}}S,\quad
  \exists a\in C_S^0,\ \operatorname{dec}_S(c)=\delta_0a\}.
\]
Equivalently, for \(S\in\mathcal S_{\mathrm{coh}}\), membership
\(c\in G_S\) is precisely the committed one-cochain coherence contract.
Consequently,
\[
  c\in G_S
  \Longrightarrow
  c\in\mathcal C_{\mathrm{com}}^{\mathrm{wf}},
\]
and the sheaf coherence hypotheses give
\[
  \delta_1\delta_0=0,
  \qquad
  \Delta\text{ is built from the adjacent coboundaries}.
\]
Moreover, exactness now applies to the particular payload:
\[
  \delta_1\operatorname{dec}_S(c)=0.
\]
Thus \(G_S\) is the accepted committed-cocycle locus. It is not the
degree-zero global-section locus. Its elements also have nonempty edge data
and a supported modulus \(q>1\).

\begin{definition}[Local syndrome target]\label{def:local-syndrome-target}
Fix sheaf coherence data \(S\). A local syndrome target is a quadruple
\[
  \mathcal L=(\Sigma,\sigma,Z,d_L)
\]
where
\[
  \sigma:\mathcal C_{\mathrm{com}}\to\Sigma,\qquad
  Z\subseteq\Sigma,\qquad
  d_L:\mathcal C_{\mathrm{com}}\to\mathbb R,
\]
and
\[
  c\in G_S\Longrightarrow \sigma(c)\in Z
  \qquad(c\in\mathcal C_{\mathrm{com}}).
\]
It determines the local code
\[
  \mathcal C_L\subseteq\mathcal C_{\mathrm{com}},\qquad
  c\in\mathcal C_L
  \Longleftrightarrow
  c\in G_S\wedge \sigma(c)\in Z.
\]
\end{definition}

The induced code-side proximity space has word carrier
\(\mathcal C_{\mathrm{com}}\), code \(\mathcal C_L\), and distance \(d_L\).
Hence every accepted committed cocycle lies in the local code.

\begin{definition}[Pl\"ucker-to-cochain projection]\label{def:plucker-cochain-projection}
Let \(P_{\mathrm{Pl}}\) be the projective Pl\"ucker quotient. A
Pl\"ucker-to-cochain projection is a map
\[
  E:P_{\mathrm{Pl}}\longrightarrow \mathcal C_{\mathrm{com}}
\]
such that, for every frame witness \(\mathcal F\),
\[
  \mathcal F\in\mathcal V_{\mathrm{fr}}
  \Longrightarrow
  E(\operatorname{Pl}(\mathcal F))\in G_S
\]
and
\[
  d_L(E(\operatorname{Pl}(\mathcal F)))\le \rho(\mathcal F).
\]
\end{definition}

A concrete payload encoder is a special case of this relation. It provides
\[
  E:P_{\mathrm{Pl}}\to \mathcal C_{\mathrm{com}},
  \qquad N_{\mathrm{edge}}\in\mathbb N,\quad N_{\mathrm{edge}}>0,
\]
and, for valid frames,
\[
  E(\operatorname{Pl}(\mathcal F))\in G_S,
  \qquad
  n_{\mathrm{edge}}(E(\operatorname{Pl}(\mathcal F)))=N_{\mathrm{edge}}.
\]
Thus the encoded edge vector has length \(N_{\mathrm{edge}}\). Since the
encoder consumes the projective coordinate, equality of Pl\"ucker coordinates
gives equality of the encoded committed cochain.

The storage-native statement is relational. A payload realization is a relation
\[
  \mathcal R\subseteq P_{\mathrm{Pl}}\times \mathcal C_{\mathrm{com}}
\]
such that
\[
  (\xi,c)\in\mathcal R
  \Longrightarrow c\in G_S,
  \qquad
  \xi=\xi'\ \text{and}\ (\xi,c)\in\mathcal R
  \Longrightarrow
  (\xi',c)\in\mathcal R.
\]
An observed frame binding consists of a hashed Pl\"ucker coordinate, an observed
proof-carrying cochain artifact, a valid frame, and a public statement proving
\[
  (\operatorname{Pl}(\mathcal F),c)\in\mathcal R.
\]
It yields
\[
  c\in\mathcal C_{\mathrm{com}}^{\mathrm{wf}},\qquad
  c\in G_S,\qquad
  c\in\mathcal C_L.
\]

\begin{proposition}[Observed Pl\"ucker binding gives a local codeword]
\label{prop:observed-binding-codeword}
Let an observed frame binding relative to \(S\) prove
\[
  (\operatorname{Pl}(\mathcal F),c)\in\mathcal R
\]
for \(\mathcal F\in\mathcal V_{\mathrm{fr}}\).
Then \(c\) is well formed, \(c\) is an accepted committed cocycle, and \(c\) is a
codeword for every local syndrome target \(\mathcal L\) over \(S\).
\end{proposition}

\begin{proof}
The payload-realization relation includes the implication
\[
  (\xi,c)\in\mathcal R
  \Longrightarrow c\in G_S.
\]
Applying it to \(\xi=\operatorname{Pl}(\mathcal F)\) gives
\(c\in G_S\), which in particular gives payload well-formedness and the
degree-one cocycle equation. For a local syndrome target
\(\mathcal L=(\Sigma,\sigma,Z,d_L)\), the target axiom gives
\[
  c\in G_S\Longrightarrow \sigma(c)\in Z.
\]
Together with \(c\in G_S\), this is exactly \(c\in\mathcal C_L\).
\end{proof}

\begin{definition}[Syndrome adequacy gate]\label{def:syndrome-adequacy-gate}
For a local syndrome target \(\mathcal L=(\Sigma,\sigma,Z,d_L)\), the syndrome-adequacy
gate is
\[
  \mathcal A_S(\mathcal L)
  \Longleftrightarrow
  \forall c\in\mathcal C_{\mathrm{com}},\quad
  c\in\mathcal C_{\mathrm{com}}^{\mathrm{wf}}\wedge(\sigma(c)\in Z)
  \Longrightarrow
  c\in G_S.
\]
\end{definition}

\begin{proposition}[An adequate syndrome gives a local codeword]\label{prop:zero-syndrome-codeword}
If \(\mathcal A_S(\mathcal L)\), \(c\in\mathcal C_{\mathrm{com}}^{\mathrm{wf}}\),
and \(\sigma(c)\in Z\), then
\[
  c\in\mathcal C_L.
\]
\end{proposition}

\begin{proof}
The adequacy implication gives \(c\in G_S\). Together with
\(\sigma(c)\in Z\), this is the defining conjunction for
\(c\in\mathcal C_L\).
\end{proof}

In the concrete commitment-sum bridge, the construction input is
\[
  c\in G_S
  \Longrightarrow
  \operatorname{sum}_{\mathrm{com}}(c)=0.
\]
The induced local syndrome target is
\[
  \Sigma=\mathbb Z,\qquad
  \sigma(c)=\operatorname{sum}_{\mathrm{com}}(c),\qquad
  Z=\{0\},
\]
with binary residual
\[
  d_L(c)=
  \begin{cases}
    0,& \operatorname{sum}_{\mathrm{com}}(c)=0,\\
    1,& \operatorname{sum}_{\mathrm{com}}(c)\ne 0.
  \end{cases}
\]
Global candidates have zero commitment-sum residual.

\subsection{Plane Extraction and RS Encoding}

Let \(I\) be the finite evaluation-index type. A two-coordinate plane word is
\[
  p:I\times\operatorname{Fin}(2)\to R.
\]
A plane extraction context is
\[
  \mathcal E=(N,\iota,a)
\]
where
\[
  N\in\mathbb N,\qquad
  \iota:I\times\operatorname{Fin}(2)\to\mathbb N,\qquad
  a:\mathbb Z\to R.
\]
It is locally nondegenerate when
\[
  \iota(i,0)\ne \iota(i,1)\qquad(i\in I).
\]
A committed cochain \(c=(q,n,\mathbf c,\tau)\) is aligned with
\(\mathcal E\) when
\[
  c\in\mathcal C_{\mathrm{com}}^{\mathrm{wf}},\qquad
  n=N,\qquad
  \iota(i,j)<|\mathbf c|\quad(i\in I,\ j\in\operatorname{Fin}(2)).
\]
The extracted plane is
\[
  p_c(i,j)=a(\mathbf c_{\iota(i,j)}).
\]

A plane target is a pair
\[
  \mathcal P=(\mathcal C_{\mathcal P},d_{\mathcal P})
\]
on plane words. A cochain-to-plane projection satisfies
\[
\begin{gathered}
  c\in\mathcal C_L\Longrightarrow c\text{ is aligned with }\mathcal E,\\
  c\in\mathcal C_L\Longrightarrow
    p_c\in\mathcal C_{\mathcal P},\\
  d_{\mathcal P}(p_c)\le d_L(c).
\end{gathered}
\]
The associated invariants are local nondegeneracy, alignment,
well-formedness, expected edge count, selected-index bounds, codeword
preservation, and distance control.

The projection also carries the naturality data used by the stage
certificates below. Write \(\mathsf r_C\) and \(\mathsf r_P\) for the
chosen restriction maps on committed cochains and plane words, respectively.
For a plane-syndrome carrier \(\Sigma_P\), it supplies maps
\(s_P:\{\text{plane words}\}\to\Sigma_P\) and
\(u:\Sigma\to\Sigma_P\) satisfying
\[
  p_{\mathsf r_C(c)}=\mathsf r_P(p_c),
  \qquad
  s_P(p_c)=u(\sigma(c)).
\]
These equations and local nondegeneracy are supplied fields of the
projection, as in the Lean definition.

Now let \(F\) be a finite field. A plane encoding is
\[
  \Phi:(\operatorname{Fin}(2)\to R)\to F.
\]
It sends a plane word \(p\) to a received word
\[
  \Phi_\ast(p):I\to F,\qquad
  \Phi_\ast(p)(i)=\Phi(p(i,\cdot)).
\]
An RS interpretation retains an independently meaningful plane target
\((\mathcal C_{\mathcal P},d_{\mathcal P})\) and supplies the two later-stage
conditions
\[
  p\in\mathcal C_{\mathcal P}
  \Longrightarrow
  \Phi_\ast(p)\in T,
  \qquad
  D(\Phi_\ast(p))\le d_{\mathcal P}(p).
\]
It also supplies a received-word restriction \(\mathsf r_T\), a syndrome
map \(s_T:F^I\to\Sigma_T\), and a map
\(v:\Sigma_P\to\Sigma_T\), with
\[
  \Phi_\ast(\mathsf r_P(p))=\mathsf r_T(\Phi_\ast(p)),
  \qquad
  s_T(\Phi_\ast(p))=v(s_P(p)).
\]
These are the later-stage restriction and syndrome naturality equations.
It therefore exposes an evaluation projection from the local syndrome space
to the RS code space, with
\[
  c\in\mathcal C_L
  \Longrightarrow
  \Phi_\ast(p_c)\in T,
  \qquad
  D(\Phi_\ast(p_c))\le d_L(c).
\]
For the canonical integer boundary used below,
\[
  F=\mathbb Z/12289\mathbb Z,\qquad
  \Phi(x_0,x_1)=x_0+g x_1,
\]
where \(g\) is the canonical RS generator defined in
\cref{sec:rs-canonical-endpoint}.

This integer encoding is not globally injective.  If the first coordinate is
shifted by \(12289\), its image in \(F\) is unchanged.  Consequently the
plane-to-RS far theorem cannot be inferred from the displayed encoding alone.
It must either restrict plane points to an admissible region on which \(\Phi\)
is faithful, or supply the required metric lower bound directly.  For a set
\(A\subseteq(\operatorname{Fin}(2)\to\mathbb Z)\), write
\[
  \operatorname{FaithfulOn}(\Phi,A)
  \quad\Longleftrightarrow\quad
  \Phi|_A\text{ is injective}.
\]
Coordinatewise membership in \(A\), together with
\(\operatorname{FaithfulOn}(\Phi,A)\), prevents two plane words from encoding
to the same received word.  This is a non-collapse precondition, not yet a
quantitative distance lower bound.

\begin{proposition}[Plane and RS handoff]\label{prop:plane-rs-handoff}
Let \(c\in\mathcal C_{\mathrm{com}}\). If \(c\in\mathcal C_L\), then
the cochain-plane projection gives an aligned plane \(p_c\), and the
plane-to-RS interpretation gives
\[
  \Phi_\ast(p_c)\in T,
  \qquad
  D(\Phi_\ast(p_c))\le d_L(c).
\]
\end{proposition}

\begin{proof}
The cochain-to-plane projection has three defining hypotheses. Applied to
\(c\in\mathcal C_L\), they give alignment of \(c\) with
\(\mathcal E\), the target-codeword assertion for \(p_c\), and
\[
  d_{\mathcal P}(p_c)\le d_L(c).
\]
The plane-to-RS interpretation maps the resulting plane codeword to an RS
codeword and gives
\[
  D(\Phi_\ast(p_c))\le d_{\mathcal P}(p_c).
\]
Composing this inequality with
\(d_{\mathcal P}(p_c)\le d_L(c)\) gives the metric conclusion.
\end{proof}

\subsection{Composed Chain, Gates, and Evidence}

A concrete projection chain is a pair
\[
  \Gamma=(\pi_{\mathrm{Pl}\to c},\pi_{c\to \mathrm{RS}})
\]
consisting of the Pl\"ucker-to-cochain projection and the plane-to-RS
interpretation. For a frame witness \(\mathcal F\), define
\[
\begin{aligned}
  c_\Gamma(\mathcal F)
    &=E(\operatorname{Pl}(\mathcal F)),\\
  p_\Gamma(\mathcal F)
    &=p_{c_\Gamma(\mathcal F)},\\
  y_\Gamma(\mathcal F)
    &=\Phi_\ast(p_\Gamma(\mathcal F)).
\end{aligned}
\]
The composed evaluation projection is
\[
  \mathcal F\longmapsto y_\Gamma(\mathcal F).
\]
Its metric handoff is the chain of inequalities
\[
  D(y_\Gamma(\mathcal F))
  \le d_{\mathcal P}(p_\Gamma(\mathcal F))
  \le d_L(c_\Gamma(\mathcal F))
  \le \rho(\mathcal F).
\]

\begin{proposition}[Composed projection chain]\label{prop:composed-chain}
If \(\mathcal F\) is a valid Pl\"ucker frame, then
\[
  y_\Gamma(\mathcal F)\in T.
\]
For every frame witness \(\mathcal F\),
\[
  D(y_\Gamma(\mathcal F))\le \rho(\mathcal F).
\]
\end{proposition}

\begin{proof}
For a valid frame, the Pl\"ucker-to-cochain projection gives
\[
  c_\Gamma(\mathcal F)\in G_S.
\]
The local syndrome target then gives \(\sigma(c_\Gamma(\mathcal F))\in Z\), so
\[
  c_\Gamma(\mathcal F)\in\mathcal C_L.
\]
The plane-to-RS interpretation maps local codewords to RS codewords; hence
\[
  y_\Gamma(\mathcal F)=\Phi_\ast(p_\Gamma(\mathcal F))\in T.
\]
For the metric statement, the plane-to-RS interpretation gives
\[
  D(y_\Gamma(\mathcal F))
  \le d_{\mathcal P}(p_\Gamma(\mathcal F)),
\]
the cochain-to-plane projection gives
\[
  d_{\mathcal P}(p_\Gamma(\mathcal F))
  \le d_L(c_\Gamma(\mathcal F)),
\]
and the Pl\"ucker-to-cochain projection gives
\[
  d_L(c_\Gamma(\mathcal F))\le \rho(\mathcal F).
\]
The desired inequality follows by transitivity.
\end{proof}

For one witness, the state bus is nested in the same order as the projection:
the Pl\"ucker/cochain state stores \(c_\Gamma(\mathcal F)\) and
\(d_L(c_\Gamma(\mathcal F))\le\rho(\mathcal F)\); the cochain/plane state
stores that predecessor together with \(p_\Gamma(\mathcal F)\) and
\(d_{\mathcal P}(p_\Gamma(\mathcal F))\le d_L(c_\Gamma(\mathcal F))\); and
the final plane/RS state stores its predecessor together with
\(y_\Gamma(\mathcal F)\) and
\(D(y_\Gamma(\mathcal F))\le d_{\mathcal P}(p_\Gamma(\mathcal F))\).
The definitional equalities record that each output is exactly the result of
its owning transition.  Transitivity yields
\[
  D(y_\Gamma(\mathcal F))\le \rho(\mathcal F).
\]

The accumulated stage certificate is ordered:
\[
  \chi_{\mathrm{Pl}\to c}
  \longrightarrow
  \chi_{c\to p}
  \longrightarrow
  \chi_{p\to \mathrm{RS}}
  \longrightarrow
  \chi_{\mathrm{stage}}.
\]
The final stage certificate is equivalent to the three explicit invariant
conditions for the Pl\"ucker/cochain, cochain/plane, and plane/RS stages.

\begin{definition}[Stage certificates]\label{def:stage-certificates}
For a chain \(\Gamma\), let
\[
  \chi_{\mathrm{Pl}\to c}(\Gamma)
\]
denote the conjunction
\[
\begin{gathered}
  \mathcal F\in\mathcal V_{\mathrm{fr}}
  \Rightarrow c_\Gamma(\mathcal F)\in G_S,\\
  \mathcal F\in\mathcal V_{\mathrm{fr}}
  \Rightarrow c_\Gamma(\mathcal F)\in\mathcal C_L,\\
  d_L(c_\Gamma(\mathcal F))\le \rho(\mathcal F).
\end{gathered}
\]
Let
\[
  \chi_{c\to p}(\Gamma)
\]
denote the conjunction of the typed restriction and syndrome naturality
squares for the cochain-to-plane transition, local nondegeneracy of
\(\mathcal E\), alignment of all local codewords, well-formedness of aligned
codewords, expected edge-count equality, selected-index bounds, plane-target
codeword preservation, and
\[
  d_{\mathcal P}(p_c)\le d_L(c).
\]
Let
\[
  \chi_{p\to\mathrm{RS}}(\Gamma)
\]
denote the conjunction
\[
\begin{gathered}
  \text{the typed restriction and syndrome naturality squares for }p\to\mathrm{RS},\\
  \Phi_\ast(p_c)\text{ is the received word associated with }c,\\
  p\in\mathcal C_{\mathcal P}\Rightarrow \Phi_\ast(p)\in T,\\
  D(\Phi_\ast(p))\le d_{\mathcal P}(p),
\end{gathered}
\]
together with the resulting composed codeword and distance consequences for
local cochains.
The accumulated stage certificate is the conjunction
\[
  \chi_{\mathrm{stage}}(\Gamma)
  =
  \chi_{\mathrm{Pl}\to c}(\Gamma)
  \wedge \chi_{c\to p}(\Gamma)
  \wedge \chi_{p\to\mathrm{RS}}(\Gamma).
\]
\end{definition}

\begin{proposition}[Stage certificate existence]\label{prop:stage-certificate-exists}
Every projection chain \(\Gamma\) satisfying the preceding definitions carries
\(\chi_{\mathrm{stage}}(\Gamma)\).
\end{proposition}

\begin{proof}
The Pl\"ucker-to-cochain projection supplies exactly the first and third clauses
of \(\chi_{\mathrm{Pl}\to c}\). The second follows because an accepted
committed cocycle is a local codeword by the defining implication
\[
  c\in G_S\Rightarrow \sigma(c)\in Z.
\]
The cochain-to-plane projection hypotheses are precisely its two typed
naturality squares, local nondegeneracy, alignment, plane-target preservation,
and the distance inequality; the well-formedness, edge-count, and
selected-index clauses are consequences of the alignment condition. The
plane-to-RS interpretation hypotheses give its two typed naturality squares,
the received-word equality, plane-codeword preservation, and the
plane-to-RS distance inequality.
Taking these three conjunctions gives \(\chi_{\mathrm{stage}}(\Gamma)\).
\end{proof}

The structural hard-gate bundle consists of the following four gates:
\begin{enumerate}[label=(\roman*)]
  \item Source metric: \(0\le\rho(\mathcal F)\), and \(\rho\) is
    Pl\"ucker-quotient invariant.
  \item Pl\"ucker metric: valid frames lie in the selected projective chart,
    and \(d_{\mathrm{Pl}}\) is quotient invariant.
  \item Quotient respect: if
    \(\operatorname{Pl}(\mathcal F)=\operatorname{Pl}(\mathcal G)\), then
    \[
    \begin{aligned}
      c_\Gamma(\mathcal F)&=c_\Gamma(\mathcal G),\\
      p_\Gamma(\mathcal F)&=p_\Gamma(\mathcal G),\\
      y_\Gamma(\mathcal F)&=y_\Gamma(\mathcal G).
    \end{aligned}
    \]
  \item Semantic alignment: valid frames give aligned cochains, and the
    plane and RS words are exactly the extraction and encoding outputs.
\end{enumerate}

Commitment acceptance has already established \(c\in G_S\) before the
projection chain begins.  Thus no converse inference from a local syndrome to
sheaf acceptance is required here.  The optional assertion
\(\sigma(c)\in Z\Rightarrow c\in G_S\) would be a separate adequacy hypothesis,
not a hard gate of this construction.

\begin{definition}[Hard-gate certificate]\label{def:hard-gate-certificate}
A hard-gate certificate for \(\Gamma\) is the tuple
\[
  \chi_{\mathrm{hard}}(\Gamma)
  =
  \bigl(
    \chi_{\mathrm{stage}}(\Gamma),\
    \gamma_{\mathrm{src}},\
    \gamma_{\mathrm{Pl}},\
    \gamma_{\mathrm{quot}},\
    \gamma_{\mathrm{align}}
  \bigr),
\]
where the four \(\gamma\)-terms are respectively the source metric gate,
Pl\"ucker metric/chart gate, Pl\"ucker-quotient respect gate, and
semantic-alignment gate.
\end{definition}

The two local restriction squares compose to a whole-chain restriction square,
and the two local syndrome squares compose to a whole-chain
coboundary/syndrome square.  The naturality certificate
\[
  \chi_{\mathrm{nat}}(\Gamma)
  =\bigl(\chi_{\mathrm{hard}}(\Gamma),\nu_\Gamma\bigr)
\]
stores the hard certificate unchanged and appends only these two composites.
It does not run either stage-local check a second time.

\begin{proposition}[Hard-gate certificate projections]\label{prop:hard-gate-projections}
A hard-gate certificate exposes the accumulated stage certificate and each of
the four structural gates. Its quotient-respect entry implies
\[
  \operatorname{Pl}(\mathcal F)=\operatorname{Pl}(\mathcal G)
  \Longrightarrow
  y_\Gamma(\mathcal F)=y_\Gamma(\mathcal G).
\]
Extending it through \(\chi_{\mathrm{nat}}(\Gamma)\) exposes the two
whole-chain naturality equations obtained by composing the stage-local
squares.
\end{proposition}

\begin{proof}
The structural statements are projections from the displayed tuple. The
quotient-respect entry contains equality of the downstream received
words whenever the Pl\"ucker coordinates agree, which is the displayed
conclusion.  For naturality, first apply the cochain-to-plane restriction
square and then the plane-to-RS restriction square; functoriality of
\(\Phi_\ast\) gives the whole-chain restriction equation.  Likewise, first
apply the plane-to-RS syndrome square and then the cochain-to-plane syndrome
square; applying the intervening syndrome map gives the whole-chain
coboundary/syndrome equation.  These are compositions of the stored local
proofs.
\end{proof}

The far-side gates are not injectivity claims. They are calibrated
non-collapse implications:
\[
  r_S\le \rho(\mathcal F)
  \Longrightarrow
  r_C\le d_L(c_\Gamma(\mathcal F)),
\]
\[
  r_C\le d_L(c_\Gamma(\mathcal F))
  \Longrightarrow
  r_P\le d_{\mathcal P}(p_\Gamma(\mathcal F)),
\]
and
\[
  r_P\le d_{\mathcal P}(p_\Gamma(\mathcal F))
  \Longrightarrow
  r_T\le D(y_\Gamma(\mathcal F)).
\]
Together they yield the far-region separation witness
\[
  r_S\le \rho(\mathcal F)
  \Longrightarrow
  r_T\le D(y_\Gamma(\mathcal F)).
\]
If two source witnesses collapse to the same final received word, the theorem
still says that any collapsed image of a far source witness is target-far. It
rules out far-to-near collapse, not all quotient collapse.

\begin{proposition}[Far-gate composition]\label{prop:far-gate-composition}
The three far-side gates imply the final far-region separation
\[
  r_S\le \rho(\mathcal F)
  \Longrightarrow
  r_T\le D(y_\Gamma(\mathcal F)).
\]
If \(y_\Gamma(\mathcal F)=y_\Gamma(\mathcal G)\) and
\(r_S\le \rho(\mathcal F)\), then
\[
  r_T\le D(y_\Gamma(\mathcal G)).
\]
\end{proposition}

\begin{proof}
Assume \(r_S\le\rho(\mathcal F)\). The first far gate gives
\[
  r_C\le d_L(c_\Gamma(\mathcal F)).
\]
The second gives
\[
  r_P\le d_{\mathcal P}(p_\Gamma(\mathcal F)),
\]
and the third gives
\[
  r_T\le D(y_\Gamma(\mathcal F)).
\]
If \(y_\Gamma(\mathcal F)=y_\Gamma(\mathcal G)\), substitute this equality in
the last inequality to obtain \(r_T\le D(y_\Gamma(\mathcal G))\).
\end{proof}

A radius calibration is
\[
  \kappa=(r^S_{\mathrm{near}},r^S_{\mathrm{far}},
          r^T_{\mathrm{near}},r^T_{\mathrm{far}})
\]
with source gap, target gap, and near compatibility. A calibrated final
certificate consists of a naturality-extended hard-gate certificate and a
far-region separation witness at the calibrated far radii. The intermediate
cochain and plane far radii remain explicit in a three-stage far pipeline, but
the final certificate records only its calibrated source and target endpoints.
It yields calibrated end-to-end readiness and hence a calibrated gap projection
into the RS code space.

\begin{definition}[Calibrated final certificate]\label{def:calibrated-final-certificate}
For a radius calibration \(\kappa\), a calibrated final certificate is
\[
  \Omega(\Gamma,\kappa)
  =
  \bigl(\chi_{\mathrm{nat}}(\Gamma),\omega_\kappa\bigr),
\]
where \(\omega_\kappa\) is the far-region separation statement
\[
  r^S_{\mathrm{far}}\le \rho(\mathcal F)
  \Longrightarrow
  r^T_{\mathrm{far}}\le D(y_\Gamma(\mathcal F)).
\]
The associated readiness condition is the conjunction
\[
  \chi_{\mathrm{stage}}(\Gamma)
  \wedge
  r^S_{\mathrm{near}}\le r^T_{\mathrm{near}}
  \wedge
  \omega_\kappa.
\]
\end{definition}

\begin{proposition}[Final certificate gives readiness]\label{prop:final-cert-ready}
A calibrated final certificate \(\Omega(\Gamma,\kappa)\) gives the readiness
condition above, and therefore gives the calibrated end-to-end gap projection
of \cref{prop:projection-pullback}.
\end{proposition}

\begin{proof}
Project \(\Omega(\Gamma,\kappa)\) onto its naturality-extended hard-gate
certificate and its far-region witness. By
\cref{prop:hard-gate-projections}, the former retains the hard-gate certificate,
which exposes \(\chi_{\mathrm{stage}}(\Gamma)\). The near compatibility
\[
  r^S_{\mathrm{near}}\le r^T_{\mathrm{near}}
\]
is part of the radius calibration \(\kappa\), and the far non-collapse
implication is \(\omega_\kappa\). These are exactly the three clauses of the
readiness condition. Combining readiness with the evaluation
projection of \cref{prop:composed-chain} gives the calibrated gap projection,
and \cref{prop:projection-pullback} proves the solver-pullback consequence.
\end{proof}

The evidence database has seven fixed slots:
\[
\begin{array}{c|l}
0 & \text{accumulated stage certificate}\\
1 & \text{source metric gate}\\
2 & \text{Pl\"ucker metric/chart gate}\\
3 & \text{Pl\"ucker-quotient respect gate}\\
4 & \text{semantic-alignment gate}\\
5 & \text{assembled restriction/syndrome naturality gate}\\
6 & \text{calibrated final far-region witness}.
\end{array}
\]
Completeness means that every required slot binds to its digest and is included
under the shared evidence root. In this paper the evidence root is a finite
index for the displayed certificate assertions: each slot names one of the
mathematical propositions above, and evidence completeness asserts that all
such proofs are present. Canonical byte encoding, domain
separation, hash binding, replay/state indexing, and live verification belong
to the subsequent protocol paper; they are not hypotheses in the mathematical
arguments below.

\begin{proposition}[Evidence completeness]\label{prop:evidence-completeness}
If the hard-gate prefix is complete, then every hard-gate slot binds to its
digest and is included under the evidence root. If the naturality slot is
complete, then the composed naturality proof binds and is included under that
root. If the final witness slot is complete, then the calibrated far-region
witness also binds and is included under the same root. A complete final
evidence certificate therefore recovers the calibrated final certificate and
calibrated end-to-end readiness.
\end{proposition}

\begin{proof}
Hard-gate completeness is the conjunction
\[
  \Bigl(\forall s\in\{0,\ldots,4\},\
    \beta_s\wedge\iota_s\Bigr)
  \wedge
  \beta_{\mathrm{root}}.
\]
Taking the corresponding conjuncts gives binding and inclusion for each hard
slot. Naturality completeness is
\[
  \beta_5\wedge\iota_5,
\]
and final-witness completeness is
\[
  \beta_6\wedge\iota_6,
\]
so the two appended-slot assertions follow. A final evidence certificate
contains the mathematical hard-gate certificate, its derived whole-chain
naturality proof, and the final far-region witness. The definition of the
calibrated final certificate adjoins the last witness to the
naturality-extended certificate, and the calibrated-readiness map is the
definition of end-to-end readiness for that final certificate.
\end{proof}

The public plane view is derived from the observed cochain artifact under the
same storage stack root. It is specified by the storage root, the cochain artifact
digest, inclusion of the cochain artifact, inclusion of the derived plane view,
alignment with \(\mathcal E\), and the equality
\[
  p=p_c.
\]
Consequently, selected plane coordinates are in bounds and the cochain digest
binds the artifact being projected.

\begin{proposition}[Public plane-view certificate]
\label{prop:public-plane-view-certificate}
A public plane-view certificate exposes
\[
\begin{gathered}
  \beta_{\mathrm{stack}},\qquad
  \iota_c,\qquad
  \iota_p,\\
  p=p_c,\qquad
  c\in\mathcal C_{\mathrm{com}}^{\mathrm{wf}},\qquad
  c\text{ is aligned with }\mathcal E,\\
  n_{\mathrm{edge}}(c)=N,\qquad
  \iota(i,j)<|\mathbf c|,\qquad
  \beta_c.
\end{gathered}
\]
\end{proposition}

\begin{proof}
The certificate is the conjunction of the storage-root binding
\(\beta_{\mathrm{stack}}\), the inclusion assertions \(\iota_c,\iota_p\), the
observed proof-carrying cochain statement, the plane word, the alignment proof,
and the equality \(p=p_c\). Taking conjuncts gives root binding, both inclusion
facts, alignment, and \(p=p_c\). The observed cochain statement contains a
well-formed proof-carrying commitment, hence
\(c\in\mathcal C_{\mathrm{com}}^{\mathrm{wf}}\). Alignment with
\(\mathcal E\) contains
\[
  n_{\mathrm{edge}}(c)=N
  \quad\text{and}\quad
  \iota(i,j)<|\mathbf c|.
\]
Finally, well-formedness of the observed proof-carrying artifact includes the
digest-binding assertion \(\beta_c\).
\end{proof}

\subsection{Finite Codes, RS Targets, and MCA}

For a finite coordinate type \(I\) and alphabet \(A\), a word is a function
\[
  w:I\to A,
\]
and a block code is a subset \(\mathcal C\subseteq A^I\). The Hamming distance is
\[
  d_H(u,v)=|\{i\in I:u_i\ne v_i\}|,
\]
and the relative distance is
\[
  \delta_H(u,v)=\frac{d_H(u,v)}{|I|}.
\]
The Hamming and relative balls around a received word \(w\) are
\[
\begin{aligned}
  B_{\mathcal C}(w,r)
    &=\{c\in \mathcal C:d_H(c,w)\le r\},\\
  B_{\mathcal C}(w,\delta)
    &=\{c\in \mathcal C:\delta_H(c,w)\le \delta\}.
\end{aligned}
\]
The list-size inequalities are
\[
  \forall w,\ |B_{\mathcal C}(w,r)|\le L,
  \qquad
  \forall w,\ |B_{\mathcal C}(w,\delta)|\le L.
\]
A count-radius closeness witness implies relative-radius closeness at
\(r/|I|\), and the relative-radius map is monotone.

For a field \(F\), an \(F\)-additive code is closed under zero, addition, and
scalar multiplication. Interleaving is defined by
\[
  \mathcal C^{\oplus T}
  =
  \{w:I\to(T\to A):\forall t\in T,\ i\mapsto w(i,t)\in \mathcal C\}.
\]
The pair word of \(f_1,f_2\in A^I\) is
\[
  (f_1,f_2)(i,0)=f_1(i),\qquad
  (f_1,f_2)(i,1)=f_2(i).
\]

An affine line of received words is
\[
  \ell(\gamma)=f_1+\gamma f_2.
\]
For \(S\subseteq I\), write
\[
  S\text{ is }\delta\text{-large}
  \Longleftrightarrow
  |S|/|I|\ge 1-\delta.
\]
The MCA bad challenges of \cite[Definition~4.3]{proximityPrize2026} are
\[
\begin{aligned}
  \operatorname{Bad}_{\mathrm{MCA}}(\mathcal C,\delta,\ell)
  =
  \{\gamma\in F:\ &\exists S\subseteq I\text{ such that}\\
  &S\text{ is }\delta\text{-large},\\
  &\ell(\gamma)|_S\in\mathcal C|_S,\\
  &(f_1,f_2)|_S\notin\mathcal C^{\oplus 2}|_S\}.
\end{aligned}
\]
Here \(\mathcal C|_S\) denotes the set of restrictions of codewords to
\(S\); the pair condition uses that same coordinate set.
The paper-facing MCA error bound is
\[
  \forall \ell,\qquad
  \frac{|\operatorname{Bad}_{\mathrm{MCA}}(\mathcal C,\delta,\ell)|}{|F|}
  \le \varepsilon.
\]
The count-based MCA condition implies the count-based correlated-agreement
condition: if too many line points are individually close, one of them is
MCA-good, and its agreement witness becomes a common agreement witness for the
line.

\begin{proposition}[MCA implies correlated agreement]\label{prop:mca-implies-ca}
Let \(\mathcal C\) be a finite \(F\)-additive block code. If \(\mathcal C\)
satisfies the count-based MCA bound at point radius \(r_p\), line radius
\(r_\ell\), and threshold \(\tau\), then \(\mathcal C\) satisfies the
count-based correlated-agreement bound at the same radii and threshold.
\end{proposition}

\begin{proof}
Fix an affine line \(\ell\). Let
\[
  B=\{\gamma\in F:\ell(\gamma)\text{ is close on some large set at radius }r_p\}
\]
and let
\[
  M=\{\gamma\in F:\gamma\text{ is not MCA-good for }\ell\}.
\]
The MCA bound says \(|M|\le\tau\). Suppose \(\tau<|B|\) but correlated
agreement fails, so no single large set witnesses line-close agreement at
radius \(r_\ell\). If \(\gamma\in B\), choose a large set \(S\) witnessing the
individual closeness of \(\ell(\gamma)\). If \(\gamma\) were MCA-good, this same
\(S\) would extend to a common line witness, contrary to the assumed failure of
correlated agreement. Hence \(\gamma\in M\). Thus \(B\subseteq M\), so
\[
  |B|\le |M|\le \tau,
\]
contradicting \(\tau<|B|\). Therefore a common large agreement set exists, which
is precisely the correlated-agreement conclusion.
\end{proof}

For Reed--Solomon, an evaluation domain is an injective map
\[
  a:I\hookrightarrow F.
\]
Parameters \((a,k)\) define the scalar RS code
\[
  \operatorname{RS}(a,k)
  =
  \{(p(a_i))_{i\in I}:\deg p<k\}.
\]
The paper alphabet \(C\subseteq(F^s)^n\) is the \(s\)-fold interleaving of the
scalar code. RS target data consist of the parameters, \(F\)-additivity, a
smooth-regime condition, list radius and list size, point and line radii, MCA
threshold, list-decoding bounds, MCA bounds, the paper Definition~4.3 MCA
bound, a received-word metric, radius normalization, and bridge conditions.
The bridge conditions specify the domain cardinality, smoothness, and
evaluation-point compatibility; restriction and coboundary/syndrome
naturality; preservation of global sections as RS codewords; equality of the
projection with the received word; metric normalization; near-distance
compatibility; far non-collapse; and radius-parameter matching.
These conditions give the calibrated hcrypt-to-RS bridge, and a correct
RS gap solver pulls back to a correct source-side gap solver.

\subsection{Canonical Reed--Solomon Endpoint}\label{sec:rs-canonical-endpoint}

The canonical endpoint fixes
\[
  F=\mathbb Z/12289\mathbb Z,\qquad n=4096,\qquad k=1024.
\]
Thus the rate is \(1/4\). The multiplicative group order is
\[
  12289-1=12288=3\cdot4096.
\]
The element \(11\in F^\times\) is fixed as a generator of order \(12288\).
Cubing it gives
\[
  g=11^3,
\]
with
\[
  \operatorname{ord}(g)=4096.
\]
The evaluation domain is
\[
  a_i=g^i,\qquad i\in\operatorname{Fin}(4096),
\]
and the order computation proves that \(i\mapsto a_i\) is injective. The
smoothness witness is the divisibility
\[
  4096\mid 12289-1.
\]

\begin{proposition}[Canonical power domain]\label{prop:canonical-power-domain}
In \(F=\mathbb Z/12289\mathbb Z\), the element \(11\) has multiplicative order
\(12288\), the element
\[
  g=11^3=1331
\]
has multiplicative order \(4096\), and the map
\[
  i\longmapsto g^i,\qquad i\in\operatorname{Fin}(4096),
\]
is injective.
\end{proposition}

\begin{proof}
The field has \(|F^\times|=12288\). Also
\[
  12288=2^{12}\cdot3,
\]
so the only prime divisors of \(12288\) are \(2\) and \(3\). Direct modular
arithmetic gives
\[
  11^{12288}\equiv1,\qquad
  11^{6144}\equiv12288=-1,\qquad
  11^{4096}\equiv6240
  \pmod{12289}.
\]
Thus \(11^{12288}=1\), while \(11^{12288/2}\ne1\) and
\(11^{12288/3}\ne1\). The standard order criterion for a finite group element
therefore gives
\[
  \operatorname{ord}(11)=12288.
\]
The order formula for powers gives
\[
  \operatorname{ord}(11^3)
  =
  \frac{\operatorname{ord}(11)}{\gcd(\operatorname{ord}(11),3)}
  =\frac{12288}{3}
  =
  4096.
\]
Now suppose \(g^i=g^j\) with
\(0\le i,j<4096\). Then
\[
  g^{i-j}=1,
\]
so \(i\equiv j\pmod{\operatorname{ord}(g)}\), hence
\[
  i\equiv j\pmod{4096}.
\]
The representatives \(i,j\) already lie in \(\{0,\ldots,4095\}\), so \(i=j\).
\end{proof}

The conservative radius below the Johnson bound is
\[
  r_J=2047,\qquad 2r_J<4096,
\]
and the radius strictly below capacity, recorded separately, is
\[
  r_{\mathrm{cap}}=3071,\qquad
  4r_{\mathrm{cap}}<3\cdot4096.
\]
The canonical real radius is
\[
  \delta_J=\frac{2047}{4096}.
\]
The canonical metric is relative Hamming distance to the RS code:
\[
  D(w)=\frac{1}{4096}\inf\{d_H(c,w):c\in \operatorname{RS}(a,1024)\}.
\]
Exact RS codewords have \(D(w)=0\), and Hamming-close words satisfy the
corresponding relative-radius bound. This supplies the radius-normalization
object for the canonical target.

\begin{proposition}[Canonical metric normalization]\label{prop:canonical-metric}
For the canonical RS target:
\[
  w\in\operatorname{RS}(a,1024)\Longrightarrow D(w)=0.
\]
If \(w\) is Hamming-close to the code at radius \(r\), then
\[
  D(w)\le \frac r{4096}.
\]
Moreover the radius conversion \(r\mapsto r/4096\) sends \(0\) to \(0\) and is
monotone.
\end{proposition}

\begin{proof}
The zero polynomial has degree \(<1024\), so the zero word is an RS codeword.
Thus the set
\[
  \{d_H(c,w):c\in\operatorname{RS}(a,1024)\}
\]
is nonempty and bounded below by \(0\). If \(w\) itself is a codeword, this set
contains \(d_H(w,w)=0\); since all Hamming distances are nonnegative, its
infimum is \(0\), and hence \(D(w)=0/4096=0\).

If \(w\) is Hamming-close at radius \(r\), then there is
\(c\in\operatorname{RS}(a,1024)\) with \(d_H(c,w)\le r\). The infimum of the
distance set is at most this displayed distance, so after dividing by \(4096\)
we obtain
\[
  D(w)\le r/4096.
\]
The map \(r\mapsto r/4096\) sends \(0\) to \(0\), and monotonicity follows from
division by the positive integer \(4096\).
\end{proof}

The canonical bridge data are
\[
  \beta=(\text{canonical smooth domain},\ \text{canonical RS target},\
  D,\ \Phi,\ \delta_J),
\]
where \(\Phi(x_0,x_1)=x_0+g x_1\). The endpoint proves the concrete
domain, smoothness, evaluation-point, metric-normalization, and radius-matching
facts. The remaining semantic theorem slots are restriction naturality,
coboundary/syndrome naturality, preservation of global sections as RS
codewords, equality of the projection with the received word, near-distance
compatibility, and far non-collapse.

\begin{definition}[Canonical bridge certificate]\label{def:canonical-bridge-certificate}
For a canonical chain \(\Gamma\), let \(\Theta_\Gamma\) be the six-tuple of
theorem slots
\[
\begin{gathered}
  \theta^{\mathrm{rest}}_\Gamma,\qquad
  \theta^{\mathrm{cob}}_\Gamma,\qquad
  \theta^{\mathrm{code}}_\Gamma,\\
  \theta^{\mathrm{word}}_\Gamma,\qquad
  \theta^{\mathrm{near}}_\Gamma,\qquad
  \theta^{\mathrm{far}}_\Gamma.
\end{gathered}
\]
A bridge certificate is a proof of their conjunction:
\[
  \mathcal B_\Gamma
  =
  \theta^{\mathrm{rest}}_\Gamma
  \wedge\theta^{\mathrm{cob}}_\Gamma
  \wedge\theta^{\mathrm{code}}_\Gamma
  \wedge\theta^{\mathrm{word}}_\Gamma
  \wedge\theta^{\mathrm{near}}_\Gamma
  \wedge\theta^{\mathrm{far}}_\Gamma.
\]
An RS metric certificate is a pair
\[
  \Omega_{\mathrm{RS}}=(\kappa,\Omega(\Gamma,\kappa)),
\]
where \(\Omega(\Gamma,\kappa)\) is the calibrated final certificate of
\cref{def:calibrated-final-certificate}.
\end{definition}

\begin{proposition}[RS metric certificate fills the bridge slots]
\label{prop:rs-metric-cert-slots}
An RS metric certificate \(\Omega_{\mathrm{RS}}\) determines a bridge
certificate for \(\Theta_\Gamma\). Consequently it completes the generic
canonical RS bridge conditions once the endpoint data
\[
  (F,n,k,a,D,\Phi,\delta_J)
\]
have been fixed as above.
\end{proposition}

\begin{proof}
Let \(\Omega_{\mathrm{RS}}=(\kappa,\Omega)\). By
\cref{prop:final-cert-ready}, \(\Omega\) exposes the accumulated stage
certificate, the derived whole-chain naturality gate, near compatibility, and
far non-collapse. The naturality gate supplies the restriction-naturality and
coboundary/syndrome-naturality slots by composing the two corresponding
stage-local squares; it does not assume new unrelated propositions. By
\cref{prop:composed-chain}, valid global frames map to RS codewords, giving
\(\theta^{\mathrm{code}}_\Gamma\). The projection-matches-word slot is the
definitional equality
\[
  \Gamma(\mathcal F)=y_\Gamma(\mathcal F).
\]
The near slot is \(\kappa\)'s near compatibility, and the far slot is
\(\Omega\)'s far-region witness. These six proofs form
\(\mathcal B_\Gamma\).

The remaining canonical bridge conditions are endpoint facts: domain
cardinality, smoothness, evaluation-point compatibility, metric normalization,
and radius matching. They were fixed in the canonical endpoint data and proved
in \cref{prop:canonical-power-domain,prop:canonical-metric}. Combining those
endpoint facts with \(\mathcal B_\Gamma\) completes the bridge conditions.
\end{proof}

A canonical Johnson far-separation witness is the statement
\[
  \delta_J\le \rho(\mathcal F)
  \Longrightarrow
  \delta_J\le D(y_\Gamma(\mathcal F)).
\]
It may be supplied directly, by a pointwise lower bound
\[
  \rho(\mathcal F)\le D(y_\Gamma(\mathcal F)),
\]
or by exact distance preservation
\[
  D(y_\Gamma(\mathcal F))=\rho(\mathcal F).
\]
Together with the naturality-extended hard-gate certificate, this witness gives
the canonical calibrated RS bridge after selecting the decision near radius.
Choose
\[
  0\le\varepsilon<\delta_J,
  \qquad
  \kappa=(\varepsilon,\delta_J,\varepsilon,\delta_J).
\]
Thus the source and target use the same pair of decision radii, with a strict
gap between near and far. For example, the canonical grid choice
\(\varepsilon=2046/4096\) is nonnegative and strictly below
\(\delta_J=2047/4096\). This choice changes only the decision near radius;
the list-decoding and MCA estimates below retain their radius \(\delta_J\).

\begin{proposition}[Canonical bridge completion]\label{prop:canonical-bridge}
Fix \(0\le\varepsilon<\delta_J\). A naturality-extended hard-gate
certificate for \(\Gamma\), together with a
canonical Johnson far-separation witness
\[
  \delta_J\le \rho(\mathcal F)
  \Longrightarrow
  \delta_J\le D(y_\Gamma(\mathcal F)),
\]
gives a completed calibrated hcrypt-to-RS bridge for the canonical endpoint
at \(\kappa=(\varepsilon,\delta_J,\varepsilon,\delta_J)\).
Both closed promise pairs are disjoint. Every correct target gap solver
pulls back along the same projection to a correct source gap solver.
\end{proposition}

\begin{proof}
The canonical endpoint supplies the concrete bridge conditions:
domain cardinality is \(|\operatorname{Fin}(4096)|=4096\), smoothness is
\(4096\mid12289-1\), evaluation points are \(a_i=g^i\) by
\cref{prop:canonical-power-domain}, the metric normalization is
\cref{prop:canonical-metric}, and the target's list, point and line radii
normalize to the relative Hamming radius \(\delta_J\). The decision radii
are separately selected as \((\varepsilon,\delta_J)\) on both sides.

The naturality-extended certificate contains the hard-gate certificate, hence
the accumulated stage certificate and structural gate bundle. Its two
naturality fields are the whole-chain composites of the restriction and
syndrome squares already stored at the cochain/plane and plane/RS stages. By
\cref{prop:composed-chain}, valid source frames
map to RS codewords and the composed map is definitionally
\(\mathcal F\mapsto y_\Gamma(\mathcal F)\). The selected calibration has
source and target gaps \(\varepsilon\le\delta_J\), and near compatibility
by reflexivity:
\[
  \varepsilon\le\varepsilon.
\]
In particular, \(\rho(\mathcal F)\le\varepsilon\) implies
\(D(y_\Gamma(\mathcal F))\le\varepsilon\) by the chain's metric handoff.
The displayed far-separation witness supplies the final far non-collapse slot.
The naturality-extended certificate and far-separation witness form the
calibrated final certificate \(\Omega(\Gamma,\kappa)\); together with
\(\kappa\), this is the RS metric certificate \(\Omega_{\mathrm{RS}}\). By
\cref{prop:rs-metric-cert-slots}, it fills the bridge theorem slots and
completes the canonical RS bridge conditions. Since
\(\varepsilon<\delta_J\), neither source nor target can be both near and
far. Applying the existing strictly separated gap-reduction view preserves
the same map and promises; \cref{prop:projection-pullback} gives the induced
pullback of correct RS gap solvers.
\end{proof}

The equal-threshold helper with all four radii equal to \(\delta_J\)
remains a conditional certificate-packaging construction. It supplies no
disjointness theorem or target solver at an attained closed boundary.
The selected strict-gap construction resolves that decision convention;
the geometric far-separation witness remains an independent hypothesis,
and no efficient solver is constructed here.

The deep list-decoding and MCA estimates are external theorem inputs
\cite{benSassonCarmonHabockKoppartySaraf2026,proximityPrize2026}.

\begin{assumption}[Cited canonical list/MCA bound]
The cited theorem boundary supplies:
\[
\begin{gathered}
  L_{\mathrm{cit}}\in\mathbb N,\qquad
  \tau_{\mathrm{cit}}\in\mathbb N,\\
  \mathcal C_{\mathrm{RS}}\text{ is }(r_J,L_{\mathrm{cit}})\text{-list decodable},\\
  \mathcal C_{\mathrm{RS}}\text{ is }(\delta_J,L_{\mathrm{cit}})\text{-list decodable},\\
  \mathcal C_{\mathrm{RS}}\text{ satisfies }(r_J,r_J,\tau_{\mathrm{cit}})\text{-MCA},\\
  \varepsilon_{\mathrm{MCA}}\le
  \frac{\tau_{\mathrm{cit}}}{|F|},
  \qquad
  \varepsilon_{\mathrm{MCA}}^{\mathrm{paper}}(\mathcal C_{\mathrm{RS}},\delta_J)
  \le
  \frac{\tau_{\mathrm{cit}}}{|F|}.
\end{gathered}
\]
Here \(\tau_{\mathrm{cit}}\) and \(L_{\mathrm{cit}}\) denote the specialized
finite witnesses. This paper does not claim
\(\tau_{\mathrm{cit}}<|F|\), or any fixed-bit security level, without a
separate numerical evaluation of the cited formula at the canonical
parameters. Thus the displayed fraction is a theorem-shaped external bound;
its nontriviality is a distinct parameter check.
\end{assumption}

\subsection{Canonical MCA Corollary}

For the canonical scalar RS endpoint, let
\[
  \mathcal W_{\mathrm{fin}}=F^{\operatorname{Fin}(4096)}.
\]
A canonical MCA line is
\[
  \ell=(f_1,f_2),\qquad
  \ell(\gamma)=f_1+\gamma f_2.
\]
The projection chain induces such a line from two source frame witnesses:
\[
  \ell_{\Gamma,\mathcal F_1,\mathcal F_2}
  =
  (y_\Gamma(\mathcal F_1),y_\Gamma(\mathcal F_2)).
\]
The canonical bad-challenge set is
\[
  \operatorname{Bad}_{\Gamma}(\mathcal F_1,\mathcal F_2)
  =
  \operatorname{Bad}_{\mathrm{MCA}}
  (\operatorname{RS}(a,1024),\delta_J,\ell_{\Gamma,\mathcal F_1,\mathcal F_2}).
\]

The paper-alphabet endpoint \(\mathcal C_{\mathrm{RS}}\subseteq(F^s)^n\) specializes here to
\(s=1\). At width one,
\[
  F\cong F^1,\qquad
  F^{4096}\cong(F^1)^{4096},
\]
and the scalar canonical bad-challenge set agrees with the paper's
\(s=1\) bad-challenge set. Therefore the imported paper Definition~4.3 MCA
bound yields the scalar canonical bound.

\begin{theorem}[Canonical RS MCA bound]\label{thm:canonical-rs-mca}
Assume the cited canonical list/MCA bound above. For the canonical \HE{}
Reed--Solomon endpoint at radius
\(\delta_J=2047/4096\), every projection-induced final affine line satisfies
\[
  \frac{|\operatorname{Bad}_{\Gamma}(\mathcal F_1,\mathcal F_2)|}{|F|}
  \le
  \frac{\tau_{\mathrm{cit}}}{|F|}.
\]
\end{theorem}

\begin{proof}
Let
\[
  \ell=\ell_{\Gamma,\mathcal F_1,\mathcal F_2}
  =
  (y_\Gamma(\mathcal F_1),y_\Gamma(\mathcal F_2)).
\]
By definition,
\[
  \operatorname{Bad}_{\Gamma}(\mathcal F_1,\mathcal F_2)
  =
  \operatorname{Bad}_{\mathrm{MCA}}
  (\operatorname{RS}(a,1024),\delta_J,\ell).
\]
It remains only to compare the scalar canonical line \(\ell\) with the paper
alphabet at width \(s=1\). The maps
\[
  F\longrightarrow F^1,\qquad a\longmapsto(\ast\mapsto a),
\]
and
\[
  F^1\longrightarrow F,\qquad u\longmapsto u(\ast),
\]
are inverse. Applying them coordinatewise gives inverse bijections between
\[
  F^{4096}
  \quad\text{and}\quad
  (F^1)^{4096}.
\]
Under this identification, scalar RS codewords are exactly width-one
paper-alphabet RS codewords: a width-one vector word is a codeword precisely
when its unique scalar component is an RS codeword. Agreement on a coordinate
set \(S\) is also preserved and reflected by the same bijection. The pair word
\((f_1,f_2)\) is preserved componentwise, so closeness of the interleaved pair
to \(C^{\oplus2}\) on \(S\) is unchanged.

Therefore the scalar bad-challenge set is equal to the paper Definition~4.3
bad-challenge set for the width-one endpoint. The cited width-one paper-MCA
theorem gives
\[
  \frac{
    |\operatorname{Bad}_{\mathrm{MCA}}
      (\operatorname{RS}(a,1024),\delta_J,\ell)|
  }{|F|}
  \le
  \frac{\tau_{\mathrm{cit}}}{|F|}.
\]
Substituting the displayed equality of bad-challenge sets gives the theorem.
\end{proof}

\begin{corollary}[Verifier-facing RS/MCA projection]\label{cor:verifier-rs-mca}
Assume the cited canonical list/MCA bound above and fix
\(0\le\varepsilon<\delta_J\). Given an accepted
hidden/geometric witness, a certified
Pl\"ucker-to-cochain-to-plane-to-RS projection chain, hard projection-gate
evidence extended through assembled naturality, and a canonical
Johnson-radius far-separation witness, the public verifier obtains a finite
Reed--Solomon received word and the strictly separated calibrated bridge at
\((\varepsilon,\delta_J)\). For any second projected witness, the resulting
affine line has canonical MCA bad-challenge fraction bounded by the cited
threshold \(\tau_{\mathrm{cit}}/|F|\).
\end{corollary}

\begin{proof}
The certified projection chain provides the maps
\[
  \mathcal F\longmapsto
  c_\Gamma(\mathcal F)\longmapsto
  p_\Gamma(\mathcal F)\longmapsto
  y_\Gamma(\mathcal F).
\]
By \cref{prop:composed-chain}, valid source data map to RS received words and
satisfy
\[
  D(y_\Gamma(\mathcal F))\le\rho(\mathcal F).
\]
The projection-gate evidence gives the structural and stage certificates and
the composed naturality equations; the Johnson far-separation witness then
gives the calibrated bridge by \cref{prop:canonical-bridge}. For any second
projected witness
\(\mathcal G\), \cref{thm:canonical-rs-mca} applied to the final affine line
\[
  (y_\Gamma(\mathcal F),y_\Gamma(\mathcal G))
\]
gives the stated bound
\[
  |\operatorname{Bad}_{\Gamma}(\mathcal F,\mathcal G)|/|F|
  \le \tau_{\mathrm{cit}}/|F|.
\]
\end{proof}

This corollary is a proximity statement. It says that accepted hidden/geometric
state can be projected to a finite public RS/MCA target with calibrated
non-collapse at the selected far region. It does not say that the projection is
injective, and it does not turn the RS/MCA endpoint into the source of HCP
hardness.

\subsection*{Lean Cross-Reference}

\begin{leanledger}
The end-to-end projection and gap vocabulary corresponds to
\lean{Proximity.Projection} and \lean{Proximity.Problem}. The Grassmannian and
Pl\"ucker source stages correspond to \lean{Proximity.Grassmannian},
\lean{Proximity.GrassmannianPlucker}, and
\lean{Proximity.GrassmannianPlucker.Concrete}. The committed-cochain and
local-syndrome stages correspond to \lean{Proximity.CochainBridge},
\lean{Proximity.PluckerCochain}, \lean{Proximity.PluckerCochain.Concrete}, and
\lean{Proximity.PluckerCochain.PublicBinding}; their accepted-cocycle predicate
is \lean{Proximity.CochainBridge.CommittedCochain.AcceptedCocycle}, which
inherits \lean{Axioms.CommittedCochainWellFormed} and
\lean{Axioms.oneCocycle_of_committedCochainWellFormed}. The plane and finite-field
encoding stages correspond to \lean{Proximity.CochainPlane.Core},
\lean{Proximity.CochainPlane.Concrete},
\lean{Proximity.CochainPlane.ReedSolomon}, and
\lean{Proximity.CochainReedSolomon}.  In particular,
\lean{CochainReedSolomon.PlaneEncoding.receivedWord_injective_on} is the
faithful-admissible-region lemma, while
\lean{CanonicalRS.integerPlaneEncoding_not_injective} records the global
modulus collision at the canonical integer boundary.

The composed chain, gates, public boundary, and evidence records correspond to
\lean{Proximity.ProjectionChain.Core}, \lean{Proximity.Gates.Stage},
\lean{Proximity.Gates.Final}, \lean{Proximity.Gates},
\lean{Proximity.GateEvidence}, and \lean{Proximity.Public}. The generic finite
code, RS, list-decoding, and CA/MCA vocabulary corresponds to
\lean{Proximity.RS.ListDecoding}, \lean{Proximity.RS.CorrelatedAgreement},
\lean{Proximity.RS.ReedSolomon}, \lean{Proximity.RS.Target}, and
\lean{Proximity.RS}. The canonical endpoint and MCA specialization correspond
to \lean{Proximity.CanonicalRS.Parameters},
\lean{Proximity.CanonicalRS.PowerDomain},
\lean{Proximity.CanonicalRS.Encoding}, \lean{Proximity.CanonicalRS.Target},
\lean{Proximity.CanonicalRS.Metric}, \lean{Proximity.CanonicalRS.BridgeData},
\lean{Proximity.CanonicalRS.BridgeWitness},
\lean{Proximity.CanonicalRS.BridgeCertificate},
\lean{Proximity.CanonicalRS.FarSeparation},
\lean{Proximity.CanonicalRS.MCA.Line}, \lean{Proximity.CanonicalRS.MCA},
\lean{Proximity.CanonicalRS.MCA.Concrete},
\lean{Proximity.CanonicalRS.CitedBounds}, and \lean{Proximity.CanonicalRS}.
The strict canonical decision convention is implemented by
\lean{CanonicalRS.BridgeData.separatedJohnsonRadiusCalibration} and
\lean{CanonicalRS.BridgeData.decisionNearRadiusReal_lt_johnsonRadiusReal}.
The original far witness is retained by
\lean{CanonicalRS.BridgeData.JohnsonRSFarSeparation.toSeparatedRSMetricCertificate};
\lean{toSeparatedBridge_radii} identifies its four decision radii, while
\lean{toSeparatedGapReduction} and \lean{pullbackSeparatedSolverCorrect}
retain the same projection and correct-solver pullback. These declarations
reuse \lean{Proximity.Application.GapBoundary}'s strict views of the existing
closed gap problems.
\end{leanledger}
